% ===========================================================================
%  Chapter 03 notes — Rings
%  Garling, A Course in Galois Theory
%
%  DIVISION OF LABOUR (see AI_INSTRUCTIONS.md §7):
%    The assistant owns the scaffold: preamble, sectioning, and the faithful
%    transcription of definitions, theorem statements, and examples.
%    The user owns everything between the SOLUTION / NOTE markers. No
%    assistant writes mathematics inside those regions in normal mode.
% ===========================================================================
\documentclass[11pt]{article}
\usepackage{coursemacros}

\title{Chapter 03 --- Rings}
\author{Mikey Ferguson}
\date{\today}

\begin{document}
\maketitle
\tableofcontents
\bigskip

% ---------------------------------------------------------------------------
\section{Overview}
% ===== NOTE 03.overview =====
% TODO(mferguson): in your own words, what is this chapter for and what does
% it let you do that the previous chapter did not?
% ===== END NOTE 03.overview =====

% ---------------------------------------------------------------------------
\section{Definitions}
% Assistant: transcribe each definition from the chapter here, in order,
% inside a definition environment, with the book's numbering in the label.

% ---------------------------------------------------------------------------
\section{Results}
% Assistant: transcribe each theorem/lemma/corollary STATEMENT here using the
% booktheorem environment, e.g. \begin{booktheorem}{03.1} ... \end{booktheorem}
% Proofs are the user's work and go in marked regions beneath each statement.

% ---------------------------------------------------------------------------
\section{Worked examples}
% ===== NOTE 03.examples =====
% TODO(mferguson): the baby examples you worked by hand while reading.
% ===== END NOTE 03.examples =====

% ---------------------------------------------------------------------------
\section{Loose ends}
% ===== NOTE 03.questions =====
% TODO(mferguson): what is still unclear, and what you want to come back to.
% ===== END NOTE 03.questions =====

\end{document}
